\chapter{The Canonical Dihedral Catalogue}
\label{chap:canonical-taxonomy}

\section{Taxonomic Ordering Principles}

The $1{,}859$ dihedral composition mosaics across $\mathbb{Z}_1$ through $\mathbb{Z}_9$ are organised according to a four-tier hierarchical classification:

\begin{enumerate}
    \item \textbf{Tier 1 (Lattice Ring Progression)}: Mosaics are stratified primarily by the underlying circular lattice order $n \in \{1, 2, \dots, 9\}$.
    
    \item \textbf{Tier 2 (Part-Count Stratification)}: Within each ring $\mathbb{Z}_n$, mosaics are stratified by the number of constituent blocks $s = \ell(\lambda) \in \{1, \dots, n\}$, corresponding to the Stirling numbers of the second kind.
    
    \item \textbf{Tier 3 (Lexicographical Shape Ordering)}: Within each stratum $s$, individual partition shapes $\lambda = (\lambda_1, \dots, \lambda_s) \vdash n$ follow standard lexicographical order:
    \begin{equation}
    (n) \prec (n-1, 1) \prec (n-2, 2) \prec \dots \prec (2, 1^{n-2}) \prec (1^n)
    \end{equation}
    
    \item \textbf{Tier 4 (Intra-Shape Ordering)}: Within each individual shape $\lambda$, non-equivalent dihedral orbits are ordered deterministically by:
    \begin{itemize}
        \item \textbf{Symmetric Block Count ($s_{\text{sym}}$ Descending)}: Fully achiral configurations precede mixed configurations, which precede chiral configurations.
        \item \textbf{Chord Energy ($E$ Descending)}: Mosaics spanning longer circular distances precede localized configurations:
        \begin{equation}
        E(\mathcal{M}) = \sum_{B \in \mathcal{M}} \sum_{u < v \in B} \min(|u - v|, n - |u - v|)^2
        \end{equation}
        \item \textbf{Canonical Representative String}: Deterministic tie-breaking is governed by the lexicographical order of the hexadecimal slash-delimited partition notation.
    \end{itemize}
\end{enumerate}

Every orbit is assigned a canonical integer identifier $1 \dots M_n$ on its respective ring $\mathbb{Z}_n$, matching its corresponding location in the visual atlas of Chapter~\ref{chap:visual-atlas}.

\section{Stratified Partition Shape Inventory}

Tables~\ref{tab:shape-census-z1-z6} through~\ref{tab:shape-census-z12} record the stratified partition shape inventory across the twelve rings $\mathbb{Z}_1$ through $\mathbb{Z}_{12}$ ($271$ integer partition shapes in total). In each table, the block-size partition shape monomial $\mathbf{z}^\lambda = \prod_{i=1}^n z_i^{c_i}$ serves as the primary row index; $\#p$ records the Fa\`a di Bruno multinomial count of set partitions of shape $\lambda \vdash n$; $\#o$ records the number of non-equivalent dihedral mosaic orbits; and $Q_\lambda(\omega) = \sum_{k \mid 2n} a_{k,\lambda} \omega^k$ is the bivariate dihedral orbit polynomial. Horizontal rules separate partition shapes of different part counts $s = \ell(\lambda) \in \{1, \dots, n\}$. Summing the stratified polynomial blocks on each ring $\mathbb{Z}_n$ recovers the aggregate orbit polynomial $\sum a_k \omega^k$ of Table~\ref{tab:census-z1-z16}.

\subsection*{Orbit Counts and Partition Derivatives}

A direct connection links the orbit polynomial $Q_\lambda(\omega) = \sum_{k \mid 2n} a_{k,\lambda} \omega^k$ to the scalar counts $\#o$ and $\#p$ through differential calculus:

\begin{enumerate}\setlength{\itemsep}{0pt}\setlength{\parskip}{0pt}
    \item \textbf{The Orbit Count Evaluation ($Q_\lambda(1) = \#o$)}:
    Setting $\omega = 1$ collapses every power $\omega^k \mapsto 1^k = 1$, leaving the unweighted sum of coefficients:
    \begin{equation}
    Q_\lambda(1) = \sum_{k \mid 2n} a_{k,\lambda} = \#o
    \end{equation}
    Evaluating at unity simply counts the total number of non-equivalent mosaic orbits of shape $\lambda$.

    \item \textbf{The Raw Partition Derivative ($Q_\lambda'(1) = \#p$)}:
    Differentiating $Q_\lambda(\omega)$ with respect to $\omega$ invokes the elementary power rule $\frac{d}{d\omega}(\omega^k) = k \omega^{k-1}$. Evaluating the resulting derivative at $\omega = 1$ yields:
    \begin{equation}
    Q_\lambda'(1) = \left. \frac{d}{d\omega} Q_\lambda(\omega) \right|_{\omega = 1} = \sum_{k \mid 2n} k \cdot a_{k,\lambda} = \#p
    \end{equation}
    Because each orbit of length $k$ comprises exactly $k$ distinct set partitions, the differentiation operator $\frac{d}{d\omega}$ mechanically weights each orbit by its cardinality $k = [D_n : \operatorname{Stab}(\mathcal{M})]$. The derivative $Q_\lambda'(1)$ thus reconstructs the full Fa\`a di Bruno multinomial count $\#p$ from the orbit polynomial.
\end{enumerate}

\noindent\textbf{Exemplar}\\[0.2em]
For the chord-pair shape $z_2^2$ on the square lattice $\mathbb{Z}_4$ (Table~\ref{tab:shape-census-z1-z6}), $Q_{z_2^2}(\omega) = 1\omega^1 + 1\omega^2$. Evaluating directly gives $Q_{z_2^2}(1) = 1(1)^1 + 1(1)^2 = 2 = \#o$. Differentiating and evaluating at $\omega = 1$:
\begin{equation}
Q_{z_2^2}'(\omega) = 1 + 2\omega \implies Q_{z_2^2}'(1) = 1 + 2(1) = 3 = \#p
\end{equation}
recovering the $3$ set partitions into two chords.

Summing across all shapes on ring $\mathbb{Z}_n$ yields the aggregate polynomial $Q_n(\omega) = \sum a_k \omega^k$ of Table~\ref{tab:census-z1-z16}. The same evaluation rules give $Q_n(1) = M_n$ for the dihedral mosaic count (OEIS A084708) and $Q_n'(1) = B_n$ for the $n$-th Bell number (OEIS A000110).

% Auto-generated by generate_extended_shape_census.cpp
\clearpage
\begingroup
\small
\setlength{\tabcolsep}{5pt}
\begin{longtable}{l r r l}
\caption{\label{tab:shape-census-z1-z6}Stratified Partition Shape Inventory for $\mathbb{Z}_1$ through $\mathbb{Z}_6$}\\
\toprule
Shape $\mathbf{z}^\lambda$ & $\#p$ & $\#o$ & Orbit Polynomial $Q_\lambda(\omega)$ \\
\midrule
\endfirsthead
\caption[]{Stratified Partition Shape Inventory for $\mathbb{Z}_1$ through $\mathbb{Z}_6$ (continued)}\\
\toprule
Shape $\mathbf{z}^\lambda$ & $\#p$ & $\#o$ & Orbit Polynomial $Q_\lambda(\omega)$ \\
\midrule
\endhead
\bottomrule
\endfoot
\midrule
\multicolumn{4}{c}{\textbf{Ring $\mathbb{Z}_{1}$}} \\
\midrule
$z_{1}$ & 1 & 1 & $1\omega^{1}$ \\
\midrule
\multicolumn{4}{c}{\textbf{Ring $\mathbb{Z}_{2}$}} \\
\midrule
$z_{2}$ & 1 & 1 & $1\omega^{1}$ \\
\arrayrulecolor{faintrule}\midrule[0.25pt]\arrayrulecolor{black}
$z_{1}^{2}$ & 1 & 1 & $1\omega^{1}$ \\
\midrule
\multicolumn{4}{c}{\textbf{Ring $\mathbb{Z}_{3}$}} \\
\midrule
$z_{3}$ & 1 & 1 & $1\omega^{1}$ \\
\arrayrulecolor{faintrule}\midrule[0.25pt]\arrayrulecolor{black}
$z_{2} z_{1}$ & 3 & 1 & $1\omega^{3}$ \\
\arrayrulecolor{faintrule}\midrule[0.25pt]\arrayrulecolor{black}
$z_{1}^{3}$ & 1 & 1 & $1\omega^{1}$ \\
\midrule
\multicolumn{4}{c}{\textbf{Ring $\mathbb{Z}_{4}$}} \\
\midrule
$z_{4}$ & 1 & 1 & $1\omega^{1}$ \\
\arrayrulecolor{faintrule}\midrule[0.25pt]\arrayrulecolor{black}
$z_{3} z_{1}$ & 4 & 1 & $1\omega^{4}$ \\
$z_{2}^{2}$ & 3 & 2 & $1\omega^{1} + 1\omega^{2}$ \\
\arrayrulecolor{faintrule}\midrule[0.25pt]\arrayrulecolor{black}
$z_{2} z_{1}^{2}$ & 6 & 2 & $1\omega^{2} + 1\omega^{4}$ \\
\arrayrulecolor{faintrule}\midrule[0.25pt]\arrayrulecolor{black}
$z_{1}^{4}$ & 1 & 1 & $1\omega^{1}$ \\
\midrule
\multicolumn{4}{c}{\textbf{Ring $\mathbb{Z}_{5}$}} \\
\midrule
$z_{5}$ & 1 & 1 & $1\omega^{1}$ \\
\arrayrulecolor{faintrule}\midrule[0.25pt]\arrayrulecolor{black}
$z_{4} z_{1}$ & 5 & 1 & $1\omega^{5}$ \\
$z_{3} z_{2}$ & 10 & 2 & $2\omega^{5}$ \\
\arrayrulecolor{faintrule}\midrule[0.25pt]\arrayrulecolor{black}
$z_{3} z_{1}^{2}$ & 10 & 2 & $2\omega^{5}$ \\
$z_{2}^{2} z_{1}$ & 15 & 3 & $3\omega^{5}$ \\
\arrayrulecolor{faintrule}\midrule[0.25pt]\arrayrulecolor{black}
$z_{2} z_{1}^{3}$ & 10 & 2 & $2\omega^{5}$ \\
\arrayrulecolor{faintrule}\midrule[0.25pt]\arrayrulecolor{black}
$z_{1}^{5}$ & 1 & 1 & $1\omega^{1}$ \\
\midrule
\multicolumn{4}{c}{\textbf{Ring $\mathbb{Z}_{6}$}} \\
\midrule
$z_{6}$ & 1 & 1 & $1\omega^{1}$ \\
\arrayrulecolor{faintrule}\midrule[0.25pt]\arrayrulecolor{black}
$z_{5} z_{1}$ & 6 & 1 & $1\omega^{6}$ \\
$z_{4} z_{2}$ & 15 & 3 & $1\omega^{3} + 2\omega^{6}$ \\
$z_{3}^{2}$ & 10 & 3 & $1\omega^{1} + 1\omega^{3} + 1\omega^{6}$ \\
\arrayrulecolor{faintrule}\midrule[0.25pt]\arrayrulecolor{black}
$z_{4} z_{1}^{2}$ & 15 & 3 & $1\omega^{3} + 2\omega^{6}$ \\
$z_{3} z_{2} z_{1}$ & 60 & 6 & $2\omega^{6} + 4\omega^{12}$ \\
$z_{2}^{3}$ & 15 & 5 & $1\omega^{1} + 1\omega^{2} + 2\omega^{3} + 1\omega^{6}$ \\
\arrayrulecolor{faintrule}\midrule[0.25pt]\arrayrulecolor{black}
$z_{3} z_{1}^{3}$ & 20 & 3 & $1\omega^{2} + 1\omega^{6} + 1\omega^{12}$ \\
$z_{2}^{2} z_{1}^{2}$ & 45 & 8 & $3\omega^{3} + 4\omega^{6} + 1\omega^{12}$ \\
\arrayrulecolor{faintrule}\midrule[0.25pt]\arrayrulecolor{black}
$z_{2} z_{1}^{4}$ & 15 & 3 & $1\omega^{3} + 2\omega^{6}$ \\
\arrayrulecolor{faintrule}\midrule[0.25pt]\arrayrulecolor{black}
$z_{1}^{6}$ & 1 & 1 & $1\omega^{1}$ \\
\bottomrule
\end{longtable}
\endgroup

\clearpage
\begingroup
\small
\setlength{\tabcolsep}{5pt}
\begin{longtable}{l r r l}
\caption{\label{tab:shape-census-z7-z8}\label{tab:shape-census-z7}\label{tab:shape-census-z8}Stratified Partition Shape Inventory for $\mathbb{Z}_7$ and $\mathbb{Z}_8$}\\
\toprule
Shape $\mathbf{z}^\lambda$ & $\#p$ & $\#o$ & Orbit Polynomial $Q_\lambda(\omega)$ \\
\midrule
\endfirsthead
\caption[]{Stratified Partition Shape Inventory for $\mathbb{Z}_7$ and $\mathbb{Z}_8$ (continued)}\\
\toprule
Shape $\mathbf{z}^\lambda$ & $\#p$ & $\#o$ & Orbit Polynomial $Q_\lambda(\omega)$ \\
\midrule
\endhead
\bottomrule
\endfoot
\midrule
\multicolumn{4}{c}{\textbf{Ring $\mathbb{Z}_{7}$}} \\
\midrule
$z_{7}$ & 1 & 1 & $1\omega^{1}$ \\
\arrayrulecolor{faintrule}\midrule[0.25pt]\arrayrulecolor{black}
$z_{6} z_{1}$ & 7 & 1 & $1\omega^{7}$ \\
$z_{5} z_{2}$ & 21 & 3 & $3\omega^{7}$ \\
$z_{4} z_{3}$ & 35 & 4 & $3\omega^{7} + 1\omega^{14}$ \\
\arrayrulecolor{faintrule}\midrule[0.25pt]\arrayrulecolor{black}
$z_{5} z_{1}^{2}$ & 21 & 3 & $3\omega^{7}$ \\
$z_{4} z_{2} z_{1}$ & 105 & 9 & $3\omega^{7} + 6\omega^{14}$ \\
$z_{3}^{2} z_{1}$ & 70 & 7 & $4\omega^{7} + 3\omega^{14}$ \\
$z_{3} z_{2}^{2}$ & 105 & 12 & $9\omega^{7} + 3\omega^{14}$ \\
\arrayrulecolor{faintrule}\midrule[0.25pt]\arrayrulecolor{black}
$z_{4} z_{1}^{3}$ & 35 & 4 & $3\omega^{7} + 1\omega^{14}$ \\
$z_{3} z_{2} z_{1}^{2}$ & 210 & 18 & $6\omega^{7} + 12\omega^{14}$ \\
$z_{2}^{3} z_{1}$ & 105 & 11 & $7\omega^{7} + 4\omega^{14}$ \\
\arrayrulecolor{faintrule}\midrule[0.25pt]\arrayrulecolor{black}
$z_{3} z_{1}^{4}$ & 35 & 4 & $3\omega^{7} + 1\omega^{14}$ \\
$z_{2}^{2} z_{1}^{3}$ & 105 & 12 & $9\omega^{7} + 3\omega^{14}$ \\
\arrayrulecolor{faintrule}\midrule[0.25pt]\arrayrulecolor{black}
$z_{2} z_{1}^{5}$ & 21 & 3 & $3\omega^{7}$ \\
\arrayrulecolor{faintrule}\midrule[0.25pt]\arrayrulecolor{black}
$z_{1}^{7}$ & 1 & 1 & $1\omega^{1}$ \\
\midrule
\multicolumn{4}{c}{\textbf{Ring $\mathbb{Z}_{8}$}} \\
\midrule
$z_{8}$ & 1 & 1 & $1\omega^{1}$ \\
\arrayrulecolor{faintrule}\midrule[0.25pt]\arrayrulecolor{black}
$z_{7} z_{1}$ & 8 & 1 & $1\omega^{8}$ \\
$z_{6} z_{2}$ & 28 & 4 & $1\omega^{4} + 3\omega^{8}$ \\
$z_{5} z_{3}$ & 56 & 5 & $3\omega^{8} + 2\omega^{16}$ \\
$z_{4}^{2}$ & 35 & 7 & $1\omega^{1} + 1\omega^{2} + 2\omega^{4} + 3\omega^{8}$ \\
\arrayrulecolor{faintrule}\midrule[0.25pt]\arrayrulecolor{black}
$z_{6} z_{1}^{2}$ & 28 & 4 & $1\omega^{4} + 3\omega^{8}$ \\
$z_{5} z_{2} z_{1}$ & 168 & 12 & $3\omega^{8} + 9\omega^{16}$ \\
$z_{4} z_{3} z_{1}$ & 280 & 19 & $3\omega^{8} + 16\omega^{16}$ \\
$z_{4} z_{2}^{2}$ & 210 & 22 & $1\omega^{2} + 4\omega^{4} + 10\omega^{8} + 7\omega^{16}$ \\
$z_{3}^{2} z_{2}$ & 280 & 25 & $2\omega^{4} + 12\omega^{8} + 11\omega^{16}$ \\
\arrayrulecolor{faintrule}\midrule[0.25pt]\arrayrulecolor{black}
$z_{5} z_{1}^{3}$ & 56 & 5 & $3\omega^{8} + 2\omega^{16}$ \\
$z_{4} z_{2} z_{1}^{2}$ & 420 & 33 & $1\omega^{4} + 12\omega^{8} + 20\omega^{16}$ \\
$z_{3}^{2} z_{1}^{2}$ & 280 & 25 & $2\omega^{4} + 12\omega^{8} + 11\omega^{16}$ \\
$z_{3} z_{2}^{2} z_{1}$ & 840 & 57 & $9\omega^{8} + 48\omega^{16}$ \\
$z_{2}^{4}$ & 105 & 17 & $1\omega^{1} + 2\omega^{2} + 5\omega^{4} + 8\omega^{8} + 1\omega^{16}$ \\
\arrayrulecolor{faintrule}\midrule[0.25pt]\arrayrulecolor{black}
$z_{4} z_{1}^{4}$ & 70 & 8 & $1\omega^{2} + 1\omega^{4} + 4\omega^{8} + 2\omega^{16}$ \\
$z_{3} z_{2} z_{1}^{3}$ & 560 & 38 & $6\omega^{8} + 32\omega^{16}$ \\
$z_{2}^{3} z_{1}^{2}$ & 420 & 39 & $3\omega^{4} + 21\omega^{8} + 15\omega^{16}$ \\
\arrayrulecolor{faintrule}\midrule[0.25pt]\arrayrulecolor{black}
$z_{3} z_{1}^{5}$ & 56 & 5 & $3\omega^{8} + 2\omega^{16}$ \\
$z_{2}^{2} z_{1}^{4}$ & 210 & 22 & $1\omega^{2} + 4\omega^{4} + 10\omega^{8} + 7\omega^{16}$ \\
\arrayrulecolor{faintrule}\midrule[0.25pt]\arrayrulecolor{black}
$z_{2} z_{1}^{6}$ & 28 & 4 & $1\omega^{4} + 3\omega^{8}$ \\
\arrayrulecolor{faintrule}\midrule[0.25pt]\arrayrulecolor{black}
$z_{1}^{8}$ & 1 & 1 & $1\omega^{1}$ \\
\bottomrule
\end{longtable}
\endgroup

\clearpage
\begingroup
\small
\setlength{\tabcolsep}{5pt}
\begin{longtable}{l r r l}
\caption{\label{tab:shape-census-z9}Stratified Partition Shape Inventory for $\mathbb{Z}_{9}$}\\
\toprule
Shape $\mathbf{z}^\lambda$ & $\#p$ & $\#o$ & Orbit Polynomial $Q_\lambda(\omega)$ \\
\midrule
\endfirsthead
\caption[]{Stratified Partition Shape Inventory for $\mathbb{Z}_{9}$ (continued)}\\
\toprule
Shape $\mathbf{z}^\lambda$ & $\#p$ & $\#o$ & Orbit Polynomial $Q_\lambda(\omega)$ \\
\midrule
\endhead
\bottomrule
\endfoot
$z_{9}$ & 1 & 1 & $1\omega^{1}$ \\
\arrayrulecolor{faintrule}\midrule[0.25pt]\arrayrulecolor{black}
$z_{8} z_{1}$ & 9 & 1 & $1\omega^{9}$ \\
$z_{7} z_{2}$ & 36 & 4 & $4\omega^{9}$ \\
$z_{6} z_{3}$ & 84 & 7 & $1\omega^{3} + 3\omega^{9} + 3\omega^{18}$ \\
$z_{5} z_{4}$ & 126 & 10 & $6\omega^{9} + 4\omega^{18}$ \\
\arrayrulecolor{faintrule}\midrule[0.25pt]\arrayrulecolor{black}
$z_{7} z_{1}^{2}$ & 36 & 4 & $4\omega^{9}$ \\
$z_{6} z_{2} z_{1}$ & 252 & 16 & $4\omega^{9} + 12\omega^{18}$ \\
$z_{5} z_{3} z_{1}$ & 504 & 28 & $28\omega^{18}$ \\
$z_{5} z_{2}^{2}$ & 378 & 30 & $18\omega^{9} + 12\omega^{18}$ \\
$z_{4}^{2} z_{1}$ & 315 & 23 & $11\omega^{9} + 12\omega^{18}$ \\
$z_{4} z_{3} z_{2}$ & 1,260 & 76 & $12\omega^{9} + 64\omega^{18}$ \\
$z_{3}^{3}$ & 280 & 25 & $1\omega^{1} + 3\omega^{3} + 12\omega^{9} + 9\omega^{18}$ \\
\arrayrulecolor{faintrule}\midrule[0.25pt]\arrayrulecolor{black}
$z_{6} z_{1}^{3}$ & 84 & 7 & $1\omega^{3} + 3\omega^{9} + 3\omega^{18}$ \\
$z_{5} z_{2} z_{1}^{2}$ & 756 & 48 & $12\omega^{9} + 36\omega^{18}$ \\
$z_{4} z_{3} z_{1}^{2}$ & 1,260 & 76 & $12\omega^{9} + 64\omega^{18}$ \\
$z_{4} z_{2}^{2} z_{1}$ & 1,890 & 114 & $18\omega^{9} + 96\omega^{18}$ \\
$z_{3}^{2} z_{2} z_{1}$ & 2,520 & 148 & $16\omega^{9} + 132\omega^{18}$ \\
$z_{3} z_{2}^{3}$ & 1,260 & 85 & $3\omega^{3} + 25\omega^{9} + 57\omega^{18}$ \\
\arrayrulecolor{faintrule}\midrule[0.25pt]\arrayrulecolor{black}
$z_{5} z_{1}^{4}$ & 126 & 10 & $6\omega^{9} + 4\omega^{18}$ \\
$z_{4} z_{2} z_{1}^{3}$ & 1,260 & 76 & $12\omega^{9} + 64\omega^{18}$ \\
$z_{3}^{2} z_{1}^{3}$ & 840 & 55 & $1\omega^{3} + 15\omega^{9} + 39\omega^{18}$ \\
$z_{3} z_{2}^{2} z_{1}^{2}$ & 3,780 & 228 & $36\omega^{9} + 192\omega^{18}$ \\
$z_{2}^{4} z_{1}$ & 945 & 65 & $25\omega^{9} + 40\omega^{18}$ \\
\arrayrulecolor{faintrule}\midrule[0.25pt]\arrayrulecolor{black}
$z_{4} z_{1}^{5}$ & 126 & 10 & $6\omega^{9} + 4\omega^{18}$ \\
$z_{3} z_{2} z_{1}^{4}$ & 1,260 & 76 & $12\omega^{9} + 64\omega^{18}$ \\
$z_{2}^{3} z_{1}^{3}$ & 1,260 & 85 & $3\omega^{3} + 25\omega^{9} + 57\omega^{18}$ \\
\arrayrulecolor{faintrule}\midrule[0.25pt]\arrayrulecolor{black}
$z_{3} z_{1}^{6}$ & 84 & 7 & $1\omega^{3} + 3\omega^{9} + 3\omega^{18}$ \\
$z_{2}^{2} z_{1}^{5}$ & 378 & 30 & $18\omega^{9} + 12\omega^{18}$ \\
\arrayrulecolor{faintrule}\midrule[0.25pt]\arrayrulecolor{black}
$z_{2} z_{1}^{7}$ & 36 & 4 & $4\omega^{9}$ \\
\arrayrulecolor{faintrule}\midrule[0.25pt]\arrayrulecolor{black}
$z_{1}^{9}$ & 1 & 1 & $1\omega^{1}$ \\
\bottomrule
\end{longtable}
\endgroup

\clearpage
\begingroup
\small
\setlength{\tabcolsep}{5pt}
\begin{longtable}{l r r l}
\caption{\label{tab:shape-census-z10}Stratified Partition Shape Inventory for $\mathbb{Z}_{10}$}\\
\toprule
Shape $\mathbf{z}^\lambda$ & $\#p$ & $\#o$ & Orbit Polynomial $Q_\lambda(\omega)$ \\
\midrule
\endfirsthead
\caption[]{Stratified Partition Shape Inventory for $\mathbb{Z}_{10}$ (continued)}\\
\toprule
Shape $\mathbf{z}^\lambda$ & $\#p$ & $\#o$ & Orbit Polynomial $Q_\lambda(\omega)$ \\
\midrule
\endhead
\bottomrule
\endfoot
$z_{10}$ & 1 & 1 & $1\omega^{1}$ \\
\arrayrulecolor{faintrule}\midrule[0.25pt]\arrayrulecolor{black}
$z_{9} z_{1}$ & 10 & 1 & $1\omega^{10}$ \\
$z_{8} z_{2}$ & 45 & 5 & $1\omega^{5} + 4\omega^{10}$ \\
$z_{7} z_{3}$ & 120 & 8 & $4\omega^{10} + 4\omega^{20}$ \\
$z_{6} z_{4}$ & 210 & 16 & $2\omega^{5} + 8\omega^{10} + 6\omega^{20}$ \\
$z_{5}^{2}$ & 126 & 13 & $1\omega^{1} + 3\omega^{5} + 7\omega^{10} + 2\omega^{20}$ \\
\arrayrulecolor{faintrule}\midrule[0.25pt]\arrayrulecolor{black}
$z_{8} z_{1}^{2}$ & 45 & 5 & $1\omega^{5} + 4\omega^{10}$ \\
$z_{7} z_{2} z_{1}$ & 360 & 20 & $4\omega^{10} + 16\omega^{20}$ \\
$z_{6} z_{3} z_{1}$ & 840 & 44 & $4\omega^{10} + 40\omega^{20}$ \\
$z_{6} z_{2}^{2}$ & 630 & 46 & $6\omega^{5} + 20\omega^{10} + 20\omega^{20}$ \\
$z_{5} z_{4} z_{1}$ & 1,260 & 66 & $6\omega^{10} + 60\omega^{20}$ \\
$z_{5} z_{3} z_{2}$ & 2,520 & 132 & $12\omega^{10} + 120\omega^{20}$ \\
$z_{4}^{2} z_{2}$ & 1,575 & 101 & $7\omega^{5} + 34\omega^{10} + 60\omega^{20}$ \\
$z_{4} z_{3}^{2}$ & 2,100 & 124 & $4\omega^{5} + 32\omega^{10} + 88\omega^{20}$ \\
\arrayrulecolor{faintrule}\midrule[0.25pt]\arrayrulecolor{black}
$z_{7} z_{1}^{3}$ & 120 & 8 & $4\omega^{10} + 4\omega^{20}$ \\
$z_{6} z_{2} z_{1}^{2}$ & 1,260 & 74 & $22\omega^{10} + 52\omega^{20}$ \\
$z_{5} z_{3} z_{1}^{2}$ & 2,520 & 132 & $12\omega^{10} + 120\omega^{20}$ \\
$z_{5} z_{2}^{2} z_{1}$ & 3,780 & 198 & $18\omega^{10} + 180\omega^{20}$ \\
$z_{4}^{2} z_{1}^{2}$ & 1,575 & 101 & $7\omega^{5} + 34\omega^{10} + 60\omega^{20}$ \\
$z_{4} z_{3} z_{2} z_{1}$ & 12,600 & 636 & $12\omega^{10} + 624\omega^{20}$ \\
$z_{4} z_{2}^{3}$ & 3,150 & 190 & $6\omega^{5} + 56\omega^{10} + 128\omega^{20}$ \\
$z_{3}^{3} z_{1}$ & 2,800 & 148 & $16\omega^{10} + 132\omega^{20}$ \\
$z_{3}^{2} z_{2}^{2}$ & 6,300 & 364 & $12\omega^{5} + 80\omega^{10} + 272\omega^{20}$ \\
\arrayrulecolor{faintrule}\midrule[0.25pt]\arrayrulecolor{black}
$z_{6} z_{1}^{4}$ & 210 & 16 & $2\omega^{5} + 8\omega^{10} + 6\omega^{20}$ \\
$z_{5} z_{2} z_{1}^{3}$ & 2,520 & 132 & $12\omega^{10} + 120\omega^{20}$ \\
$z_{4} z_{3} z_{1}^{3}$ & 4,200 & 216 & $12\omega^{10} + 204\omega^{20}$ \\
$z_{4} z_{2}^{2} z_{1}^{2}$ & 9,450 & 516 & $6\omega^{5} + 78\omega^{10} + 432\omega^{20}$ \\
$z_{3}^{2} z_{2} z_{1}^{2}$ & 12,600 & 668 & $76\omega^{10} + 592\omega^{20}$ \\
$z_{3} z_{2}^{3} z_{1}$ & 12,600 & 644 & $28\omega^{10} + 616\omega^{20}$ \\
$z_{2}^{5}$ & 945 & 79 & $1\omega^{1} + 2\omega^{2} + 12\omega^{5} + 40\omega^{10} + 24\omega^{20}$ \\
\arrayrulecolor{faintrule}\midrule[0.25pt]\arrayrulecolor{black}
$z_{5} z_{1}^{5}$ & 252 & 16 & $1\omega^{2} + 5\omega^{10} + 10\omega^{20}$ \\
$z_{4} z_{2} z_{1}^{4}$ & 3,150 & 174 & $2\omega^{5} + 30\omega^{10} + 142\omega^{20}$ \\
$z_{3}^{2} z_{1}^{4}$ & 2,100 & 124 & $4\omega^{5} + 32\omega^{10} + 88\omega^{20}$ \\
$z_{3} z_{2}^{2} z_{1}^{3}$ & 12,600 & 648 & $36\omega^{10} + 612\omega^{20}$ \\
$z_{2}^{4} z_{1}^{2}$ & 4,725 & 287 & $13\omega^{5} + 82\omega^{10} + 192\omega^{20}$ \\
\arrayrulecolor{faintrule}\midrule[0.25pt]\arrayrulecolor{black}
$z_{4} z_{1}^{6}$ & 210 & 16 & $2\omega^{5} + 8\omega^{10} + 6\omega^{20}$ \\
$z_{3} z_{2} z_{1}^{5}$ & 2,520 & 132 & $12\omega^{10} + 120\omega^{20}$ \\
$z_{2}^{3} z_{1}^{4}$ & 3,150 & 190 & $6\omega^{5} + 56\omega^{10} + 128\omega^{20}$ \\
\arrayrulecolor{faintrule}\midrule[0.25pt]\arrayrulecolor{black}
$z_{3} z_{1}^{7}$ & 120 & 8 & $4\omega^{10} + 4\omega^{20}$ \\
$z_{2}^{2} z_{1}^{6}$ & 630 & 46 & $6\omega^{5} + 20\omega^{10} + 20\omega^{20}$ \\
\arrayrulecolor{faintrule}\midrule[0.25pt]\arrayrulecolor{black}
$z_{2} z_{1}^{8}$ & 45 & 5 & $1\omega^{5} + 4\omega^{10}$ \\
\arrayrulecolor{faintrule}\midrule[0.25pt]\arrayrulecolor{black}
$z_{1}^{10}$ & 1 & 1 & $1\omega^{1}$ \\
\bottomrule
\end{longtable}
\endgroup

\clearpage
\begingroup
\small
\setlength{\tabcolsep}{5pt}
\begin{longtable}{l r r l}
\caption{\label{tab:shape-census-z11}Stratified Partition Shape Inventory for $\mathbb{Z}_{11}$}\\
\toprule
Shape $\mathbf{z}^\lambda$ & $\#p$ & $\#o$ & Orbit Polynomial $Q_\lambda(\omega)$ \\
\midrule
\endfirsthead
\caption[]{Stratified Partition Shape Inventory for $\mathbb{Z}_{11}$ (continued)}\\
\toprule
Shape $\mathbf{z}^\lambda$ & $\#p$ & $\#o$ & Orbit Polynomial $Q_\lambda(\omega)$ \\
\midrule
\endhead
\bottomrule
\endfoot
$z_{11}$ & 1 & 1 & $1\omega^{1}$ \\
\arrayrulecolor{faintrule}\midrule[0.25pt]\arrayrulecolor{black}
$z_{10} z_{1}$ & 11 & 1 & $1\omega^{11}$ \\
$z_{9} z_{2}$ & 55 & 5 & $5\omega^{11}$ \\
$z_{8} z_{3}$ & 165 & 10 & $5\omega^{11} + 5\omega^{22}$ \\
$z_{7} z_{4}$ & 330 & 20 & $10\omega^{11} + 10\omega^{22}$ \\
$z_{6} z_{5}$ & 462 & 26 & $10\omega^{11} + 16\omega^{22}$ \\
\arrayrulecolor{faintrule}\midrule[0.25pt]\arrayrulecolor{black}
$z_{9} z_{1}^{2}$ & 55 & 5 & $5\omega^{11}$ \\
$z_{8} z_{2} z_{1}$ & 495 & 25 & $5\omega^{11} + 20\omega^{22}$ \\
$z_{7} z_{3} z_{1}$ & 1,320 & 60 & $60\omega^{22}$ \\
$z_{7} z_{2}^{2}$ & 990 & 60 & $30\omega^{11} + 30\omega^{22}$ \\
$z_{6} z_{4} z_{1}$ & 2,310 & 110 & $10\omega^{11} + 100\omega^{22}$ \\
$z_{6} z_{3} z_{2}$ & 4,620 & 220 & $20\omega^{11} + 200\omega^{22}$ \\
$z_{5}^{2} z_{1}$ & 1,386 & 71 & $16\omega^{11} + 55\omega^{22}$ \\
$z_{5} z_{4} z_{2}$ & 6,930 & 330 & $30\omega^{11} + 300\omega^{22}$ \\
$z_{5} z_{3}^{2}$ & 4,620 & 230 & $40\omega^{11} + 190\omega^{22}$ \\
$z_{4}^{2} z_{3}$ & 5,775 & 290 & $55\omega^{11} + 235\omega^{22}$ \\
\arrayrulecolor{faintrule}\midrule[0.25pt]\arrayrulecolor{black}
$z_{8} z_{1}^{3}$ & 165 & 10 & $5\omega^{11} + 5\omega^{22}$ \\
$z_{7} z_{2} z_{1}^{2}$ & 1,980 & 100 & $20\omega^{11} + 80\omega^{22}$ \\
$z_{6} z_{3} z_{1}^{2}$ & 4,620 & 220 & $20\omega^{11} + 200\omega^{22}$ \\
$z_{6} z_{2}^{2} z_{1}$ & 6,930 & 330 & $30\omega^{11} + 300\omega^{22}$ \\
$z_{5} z_{4} z_{1}^{2}$ & 6,930 & 330 & $30\omega^{11} + 300\omega^{22}$ \\
$z_{5} z_{3} z_{2} z_{1}$ & 27,720 & 1,260 & $1,260\omega^{22}$ \\
$z_{5} z_{2}^{3}$ & 6,930 & 350 & $70\omega^{11} + 280\omega^{22}$ \\
$z_{4}^{2} z_{2} z_{1}$ & 17,325 & 815 & $55\omega^{11} + 760\omega^{22}$ \\
$z_{4} z_{3}^{2} z_{1}$ & 23,100 & 1,070 & $40\omega^{11} + 1,030\omega^{22}$ \\
$z_{4} z_{3} z_{2}^{2}$ & 34,650 & 1,620 & $90\omega^{11} + 1,530\omega^{22}$ \\
$z_{3}^{3} z_{2}$ & 15,400 & 740 & $80\omega^{11} + 660\omega^{22}$ \\
\arrayrulecolor{faintrule}\midrule[0.25pt]\arrayrulecolor{black}
$z_{7} z_{1}^{4}$ & 330 & 20 & $10\omega^{11} + 10\omega^{22}$ \\
$z_{6} z_{2} z_{1}^{3}$ & 4,620 & 220 & $20\omega^{11} + 200\omega^{22}$ \\
$z_{5} z_{3} z_{1}^{3}$ & 9,240 & 420 & $420\omega^{22}$ \\
$z_{5} z_{2}^{2} z_{1}^{2}$ & 20,790 & 990 & $90\omega^{11} + 900\omega^{22}$ \\
$z_{4}^{2} z_{1}^{3}$ & 5,775 & 290 & $55\omega^{11} + 235\omega^{22}$ \\
$z_{4} z_{3} z_{2} z_{1}^{2}$ & 69,300 & 3,180 & $60\omega^{11} + 3,120\omega^{22}$ \\
$z_{4} z_{2}^{3} z_{1}$ & 34,650 & 1,610 & $70\omega^{11} + 1,540\omega^{22}$ \\
$z_{3}^{3} z_{1}^{2}$ & 15,400 & 740 & $80\omega^{11} + 660\omega^{22}$ \\
$z_{3}^{2} z_{2}^{2} z_{1}$ & 69,300 & 3,210 & $120\omega^{11} + 3,090\omega^{22}$ \\
$z_{3} z_{2}^{4}$ & 17,325 & 850 & $125\omega^{11} + 725\omega^{22}$ \\
\arrayrulecolor{faintrule}\midrule[0.25pt]\arrayrulecolor{black}
$z_{6} z_{1}^{5}$ & 462 & 26 & $10\omega^{11} + 16\omega^{22}$ \\
$z_{5} z_{2} z_{1}^{4}$ & 6,930 & 330 & $30\omega^{11} + 300\omega^{22}$ \\
$z_{4} z_{3} z_{1}^{4}$ & 11,550 & 540 & $30\omega^{11} + 510\omega^{22}$ \\
$z_{4} z_{2}^{2} z_{1}^{3}$ & 34,650 & 1,620 & $90\omega^{11} + 1,530\omega^{22}$ \\
$z_{3}^{2} z_{2} z_{1}^{3}$ & 46,200 & 2,140 & $80\omega^{11} + 2,060\omega^{22}$ \\
$z_{3} z_{2}^{3} z_{1}^{2}$ & 69,300 & 3,220 & $140\omega^{11} + 3,080\omega^{22}$ \\
$z_{2}^{5} z_{1}$ & 10,395 & 513 & $81\omega^{11} + 432\omega^{22}$ \\
\arrayrulecolor{faintrule}\midrule[0.25pt]\arrayrulecolor{black}
$z_{5} z_{1}^{6}$ & 462 & 26 & $10\omega^{11} + 16\omega^{22}$ \\
$z_{4} z_{2} z_{1}^{5}$ & 6,930 & 330 & $30\omega^{11} + 300\omega^{22}$ \\
$z_{3}^{2} z_{1}^{5}$ & 4,620 & 230 & $40\omega^{11} + 190\omega^{22}$ \\
$z_{3} z_{2}^{2} z_{1}^{4}$ & 34,650 & 1,620 & $90\omega^{11} + 1,530\omega^{22}$ \\
$z_{2}^{4} z_{1}^{3}$ & 17,325 & 850 & $125\omega^{11} + 725\omega^{22}$ \\
\arrayrulecolor{faintrule}\midrule[0.25pt]\arrayrulecolor{black}
$z_{4} z_{1}^{7}$ & 330 & 20 & $10\omega^{11} + 10\omega^{22}$ \\
$z_{3} z_{2} z_{1}^{6}$ & 4,620 & 220 & $20\omega^{11} + 200\omega^{22}$ \\
$z_{2}^{3} z_{1}^{5}$ & 6,930 & 350 & $70\omega^{11} + 280\omega^{22}$ \\
\arrayrulecolor{faintrule}\midrule[0.25pt]\arrayrulecolor{black}
$z_{3} z_{1}^{8}$ & 165 & 10 & $5\omega^{11} + 5\omega^{22}$ \\
$z_{2}^{2} z_{1}^{7}$ & 990 & 60 & $30\omega^{11} + 30\omega^{22}$ \\
\arrayrulecolor{faintrule}\midrule[0.25pt]\arrayrulecolor{black}
$z_{2} z_{1}^{9}$ & 55 & 5 & $5\omega^{11}$ \\
\arrayrulecolor{faintrule}\midrule[0.25pt]\arrayrulecolor{black}
$z_{1}^{11}$ & 1 & 1 & $1\omega^{1}$ \\
\bottomrule
\end{longtable}
\endgroup

\clearpage
\begingroup
\small
\setlength{\tabcolsep}{5pt}
\begin{longtable}{l r r l}
\caption{\label{tab:shape-census-z12}\label{tab:shape-census-z1-z12}Stratified Partition Shape Inventory for $\mathbb{Z}_{12}$}\\
\toprule
Shape $\mathbf{z}^\lambda$ & $\#p$ & $\#o$ & Orbit Polynomial $Q_\lambda(\omega)$ \\
\midrule
\endfirsthead
\caption[]{Stratified Partition Shape Inventory for $\mathbb{Z}_{12}$ (continued)}\\
\toprule
Shape $\mathbf{z}^\lambda$ & $\#p$ & $\#o$ & Orbit Polynomial $Q_\lambda(\omega)$ \\
\midrule
\endhead
\bottomrule
\endfoot
$z_{12}$ & 1 & 1 & $1\omega^{1}$ \\
\arrayrulecolor{faintrule}\midrule[0.25pt]\arrayrulecolor{black}
$z_{11} z_{1}$ & 12 & 1 & $1\omega^{12}$ \\
$z_{10} z_{2}$ & 66 & 6 & $1\omega^{6} + 5\omega^{12}$ \\
$z_{9} z_{3}$ & 220 & 12 & $1\omega^{4} + 4\omega^{12} + 7\omega^{24}$ \\
$z_{8} z_{4}$ & 495 & 29 & $1\omega^{3} + 2\omega^{6} + 12\omega^{12} + 14\omega^{24}$ \\
$z_{7} z_{5}$ & 792 & 38 & $10\omega^{12} + 28\omega^{24}$ \\
$z_{6}^{2}$ & 462 & 35 & \begin{tabular}[t]{@{}l@{}} $1\omega^{1} + 1\omega^{2} + 1\omega^{3} + 4\omega^{6} + {}$ \\ $\quad 20\omega^{12} + 8\omega^{24}$ \end{tabular} \\
\arrayrulecolor{faintrule}\midrule[0.25pt]\arrayrulecolor{black}
$z_{10} z_{1}^{2}$ & 66 & 6 & $1\omega^{6} + 5\omega^{12}$ \\
$z_{9} z_{2} z_{1}$ & 660 & 30 & $5\omega^{12} + 25\omega^{24}$ \\
$z_{8} z_{3} z_{1}$ & 1,980 & 85 & $5\omega^{12} + 80\omega^{24}$ \\
$z_{8} z_{2}^{2}$ & 1,485 & 84 & $1\omega^{3} + 7\omega^{6} + 32\omega^{12} + 44\omega^{24}$ \\
$z_{7} z_{4} z_{1}$ & 3,960 & 170 & $10\omega^{12} + 160\omega^{24}$ \\
$z_{7} z_{3} z_{2}$ & 7,920 & 340 & $20\omega^{12} + 320\omega^{24}$ \\
$z_{6} z_{5} z_{1}$ & 5,544 & 236 & $10\omega^{12} + 226\omega^{24}$ \\
$z_{6} z_{4} z_{2}$ & 13,860 & 610 & $2\omega^{6} + 62\omega^{12} + 546\omega^{24}$ \\
$z_{6} z_{3}^{2}$ & 9,240 & 424 & $1\omega^{2} + 1\omega^{4} + 3\omega^{6} + 70\omega^{12} + 349\omega^{24}$ \\
$z_{5}^{2} z_{2}$ & 8,316 & 386 & $4\omega^{6} + 73\omega^{12} + 309\omega^{24}$ \\
$z_{5} z_{4} z_{3}$ & 27,720 & 1,170 & $30\omega^{12} + 1,140\omega^{24}$ \\
$z_{4}^{3}$ & 5,775 & 297 & \begin{tabular}[t]{@{}l@{}} $1\omega^{1} + 1\omega^{2} + 2\omega^{3} + 4\omega^{4} + {}$ \\ $\quad 9\omega^{6} + 1\omega^{8} + 84\omega^{12} + 195\omega^{24}$ \end{tabular} \\
\arrayrulecolor{faintrule}\midrule[0.25pt]\arrayrulecolor{black}
$z_{9} z_{1}^{3}$ & 220 & 12 & $1\omega^{4} + 4\omega^{12} + 7\omega^{24}$ \\
$z_{8} z_{2} z_{1}^{2}$ & 2,970 & 140 & $1\omega^{6} + 31\omega^{12} + 108\omega^{24}$ \\
$z_{7} z_{3} z_{1}^{2}$ & 7,920 & 340 & $20\omega^{12} + 320\omega^{24}$ \\
$z_{7} z_{2}^{2} z_{1}$ & 11,880 & 510 & $30\omega^{12} + 480\omega^{24}$ \\
$z_{6} z_{4} z_{1}^{2}$ & 13,860 & 610 & $2\omega^{6} + 62\omega^{12} + 546\omega^{24}$ \\
$z_{6} z_{3} z_{2} z_{1}$ & 55,440 & 2,320 & $20\omega^{12} + 2,300\omega^{24}$ \\
$z_{6} z_{2}^{3}$ & 13,860 & 645 & $1\omega^{2} + 4\omega^{4} + 5\omega^{6} + 119\omega^{12} + 516\omega^{24}$ \\
$z_{5}^{2} z_{1}^{2}$ & 8,316 & 386 & $4\omega^{6} + 73\omega^{12} + 309\omega^{24}$ \\
$z_{5} z_{4} z_{2} z_{1}$ & 83,160 & 3,480 & $30\omega^{12} + 3,450\omega^{24}$ \\
$z_{5} z_{3}^{2} z_{1}$ & 55,440 & 2,330 & $40\omega^{12} + 2,290\omega^{24}$ \\
$z_{5} z_{3} z_{2}^{2}$ & 83,160 & 3,510 & $90\omega^{12} + 3,420\omega^{24}$ \\
$z_{4}^{2} z_{3} z_{1}$ & 69,300 & 2,915 & $55\omega^{12} + 2,860\omega^{24}$ \\
$z_{4}^{2} z_{2}^{2}$ & 51,975 & 2,347 & $3\omega^{3} + 35\omega^{6} + 305\omega^{12} + 2,004\omega^{24}$ \\
$z_{4} z_{3}^{2} z_{2}$ & 138,600 & 5,890 & $4\omega^{6} + 224\omega^{12} + 5,662\omega^{24}$ \\
$z_{3}^{4}$ & 15,400 & 713 & \begin{tabular}[t]{@{}l@{}} $1\omega^{1} + 3\omega^{3} + 3\omega^{4} + 9\omega^{6} + {}$ \\ $\quad 3\omega^{8} + 113\omega^{12} + 581\omega^{24}$ \end{tabular} \\
\arrayrulecolor{faintrule}\midrule[0.25pt]\arrayrulecolor{black}
$z_{8} z_{1}^{4}$ & 495 & 29 & $1\omega^{3} + 2\omega^{6} + 12\omega^{12} + 14\omega^{24}$ \\
$z_{7} z_{2} z_{1}^{3}$ & 7,920 & 340 & $20\omega^{12} + 320\omega^{24}$ \\
$z_{6} z_{3} z_{1}^{3}$ & 18,480 & 781 & $1\omega^{4} + 1\omega^{8} + 19\omega^{12} + 760\omega^{24}$ \\
$z_{6} z_{2}^{2} z_{1}^{2}$ & 41,580 & 1,820 & $6\omega^{6} + 166\omega^{12} + 1,648\omega^{24}$ \\
$z_{5} z_{4} z_{1}^{3}$ & 27,720 & 1,170 & $30\omega^{12} + 1,140\omega^{24}$ \\
$z_{5} z_{3} z_{2} z_{1}^{2}$ & 166,320 & 6,960 & $60\omega^{12} + 6,900\omega^{24}$ \\
$z_{5} z_{2}^{3} z_{1}$ & 83,160 & 3,500 & $70\omega^{12} + 3,430\omega^{24}$ \\
$z_{4}^{2} z_{2} z_{1}^{2}$ & 103,950 & 4,470 & $7\omega^{6} + 267\omega^{12} + 4,196\omega^{24}$ \\
$z_{4} z_{3}^{2} z_{1}^{2}$ & 138,600 & 5,890 & $4\omega^{6} + 224\omega^{12} + 5,662\omega^{24}$ \\
$z_{4} z_{3} z_{2}^{2} z_{1}$ & 415,800 & 17,370 & $90\omega^{12} + 17,280\omega^{24}$ \\
$z_{4} z_{2}^{4}$ & 51,975 & 2,325 & $5\omega^{3} + 24\omega^{6} + 274\omega^{12} + 2,022\omega^{24}$ \\
$z_{3}^{3} z_{2} z_{1}$ & 184,800 & 7,740 & $80\omega^{12} + 7,660\omega^{24}$ \\
$z_{3}^{2} z_{2}^{3}$ & 138,600 & 6,005 & $1\omega^{2} + 4\omega^{4} + 11\omega^{6} + 435\omega^{12} + 5,554\omega^{24}$ \\
\arrayrulecolor{faintrule}\midrule[0.25pt]\arrayrulecolor{black}
$z_{7} z_{1}^{5}$ & 792 & 38 & $10\omega^{12} + 28\omega^{24}$ \\
$z_{6} z_{2} z_{1}^{4}$ & 13,860 & 610 & $2\omega^{6} + 62\omega^{12} + 546\omega^{24}$ \\
$z_{5} z_{3} z_{1}^{4}$ & 27,720 & 1,170 & $30\omega^{12} + 1,140\omega^{24}$ \\
$z_{5} z_{2}^{2} z_{1}^{3}$ & 83,160 & 3,510 & $90\omega^{12} + 3,420\omega^{24}$ \\
$z_{4}^{2} z_{1}^{4}$ & 17,325 & 792 & $3\omega^{3} + 12\omega^{6} + 117\omega^{12} + 660\omega^{24}$ \\
$z_{4} z_{3} z_{2} z_{1}^{3}$ & 277,200 & 11,580 & $60\omega^{12} + 11,520\omega^{24}$ \\
$z_{4} z_{2}^{3} z_{1}^{2}$ & 207,900 & 8,860 & $6\omega^{6} + 386\omega^{12} + 8,468\omega^{24}$ \\
$z_{3}^{3} z_{1}^{3}$ & 61,600 & 2,610 & $4\omega^{4} + 3\omega^{8} + 76\omega^{12} + 2,527\omega^{24}$ \\
$z_{3}^{2} z_{2}^{2} z_{1}^{2}$ & 415,800 & 17,630 & $12\omega^{6} + 592\omega^{12} + 17,026\omega^{24}$ \\
$z_{3} z_{2}^{4} z_{1}$ & 207,900 & 8,725 & $125\omega^{12} + 8,600\omega^{24}$ \\
$z_{2}^{6}$ & 10,395 & 554 & \begin{tabular}[t]{@{}l@{}} $1\omega^{1} + 3\omega^{2} + 4\omega^{3} + 5\omega^{4} + {}$ \\ $\quad 30\omega^{6} + 174\omega^{12} + 337\omega^{24}$ \end{tabular} \\
\arrayrulecolor{faintrule}\midrule[0.25pt]\arrayrulecolor{black}
$z_{6} z_{1}^{6}$ & 924 & 50 & $1\omega^{2} + 1\omega^{4} + 1\omega^{6} + 18\omega^{12} + 29\omega^{24}$ \\
$z_{5} z_{2} z_{1}^{5}$ & 16,632 & 708 & $30\omega^{12} + 678\omega^{24}$ \\
$z_{4} z_{3} z_{1}^{5}$ & 27,720 & 1,170 & $30\omega^{12} + 1,140\omega^{24}$ \\
$z_{4} z_{2}^{2} z_{1}^{4}$ & 103,950 & 4,463 & $17\omega^{6} + 238\omega^{12} + 4,208\omega^{24}$ \\
$z_{3}^{2} z_{2} z_{1}^{4}$ & 138,600 & 5,890 & $4\omega^{6} + 224\omega^{12} + 5,662\omega^{24}$ \\
$z_{3} z_{2}^{3} z_{1}^{3}$ & 277,200 & 11,623 & $3\omega^{4} + 3\omega^{8} + 137\omega^{12} + 11,480\omega^{24}$ \\
$z_{2}^{5} z_{1}^{2}$ & 62,370 & 2,792 & $13\omega^{6} + 367\omega^{12} + 2,412\omega^{24}$ \\
\arrayrulecolor{faintrule}\midrule[0.25pt]\arrayrulecolor{black}
$z_{5} z_{1}^{7}$ & 792 & 38 & $10\omega^{12} + 28\omega^{24}$ \\
$z_{4} z_{2} z_{1}^{6}$ & 13,860 & 610 & $2\omega^{6} + 62\omega^{12} + 546\omega^{24}$ \\
$z_{3}^{2} z_{1}^{6}$ & 9,240 & 424 & $1\omega^{2} + 1\omega^{4} + 3\omega^{6} + 70\omega^{12} + 349\omega^{24}$ \\
$z_{3} z_{2}^{2} z_{1}^{5}$ & 83,160 & 3,510 & $90\omega^{12} + 3,420\omega^{24}$ \\
$z_{2}^{4} z_{1}^{4}$ & 51,975 & 2,325 & $5\omega^{3} + 24\omega^{6} + 274\omega^{12} + 2,022\omega^{24}$ \\
\arrayrulecolor{faintrule}\midrule[0.25pt]\arrayrulecolor{black}
$z_{4} z_{1}^{8}$ & 495 & 29 & $1\omega^{3} + 2\omega^{6} + 12\omega^{12} + 14\omega^{24}$ \\
$z_{3} z_{2} z_{1}^{7}$ & 7,920 & 340 & $20\omega^{12} + 320\omega^{24}$ \\
$z_{2}^{3} z_{1}^{6}$ & 13,860 & 645 & $1\omega^{2} + 4\omega^{4} + 5\omega^{6} + 119\omega^{12} + 516\omega^{24}$ \\
\arrayrulecolor{faintrule}\midrule[0.25pt]\arrayrulecolor{black}
$z_{3} z_{1}^{9}$ & 220 & 12 & $1\omega^{4} + 4\omega^{12} + 7\omega^{24}$ \\
$z_{2}^{2} z_{1}^{8}$ & 1,485 & 84 & $1\omega^{3} + 7\omega^{6} + 32\omega^{12} + 44\omega^{24}$ \\
\arrayrulecolor{faintrule}\midrule[0.25pt]\arrayrulecolor{black}
$z_{2} z_{1}^{10}$ & 66 & 6 & $1\omega^{6} + 5\omega^{12}$ \\
\arrayrulecolor{faintrule}\midrule[0.25pt]\arrayrulecolor{black}
$z_{1}^{12}$ & 1 & 1 & $1\omega^{1}$ \\
\bottomrule
\end{longtable}
\endgroup



\section{The Relational Datasets and Embedding Manifest}
\label{sec:relational-datasets}

The complete analytical inventory across all $219{,}467$ dihedral mosaics on $\mathbb{Z}_1$ through $\mathbb{Z}_{12}$ is structured into two normalised, greppable tables housed in \path{data/} (provided in tab-separated \texttt{.tsv} and comma-separated \texttt{.csv} formats):

\begin{enumerate}
    \item \textbf{Master Mosaic Catalogue (\path{data/mosaics.tsv}, \path{data/mosaics.csv})}:
    Contains all $219{,}467$ records across $\mathbb{Z}_1..\mathbb{Z}_{12}$ ($509$ records on $\mathbb{Z}_1..\mathbb{Z}_8$, $1{,}350$ on $\mathbb{Z}_9$, $6{,}351$ on $\mathbb{Z}_{10}$, $31{,}950$ on $\mathbb{Z}_{11}$, and $179{,}307$ on $\mathbb{Z}_{12}$). Fields include ring order $n$, ring-local ID ($1 \dots M_n$), shape $\lambda$, canonical slash representation, dihedral invariants ($|\mathcal{O}|$, $|\operatorname{Stab}|$, $s_{\text{sym}}$, chirality signature), circular chord energy $E$, polyatlas polygon profile, and the relational \emph{block manifest}. The first $1{,}859$ records correspond directly to the visual plates in Chapter~\ref{chap:visual-atlas}.
    
    \item \textbf{Relational Block Bridge (\path{data/mosaic_blocks.tsv}, \path{data/mosaic_blocks.csv})}:
    Contains all constituent blocks across the $219{,}467$ canonical mosaics, linking every block directly into the single inscribed polygon census of \path{polyatlas/data/inngons.tsv}.
\end{enumerate}

\noindent\textbf{The Block Embedding Manifest}\\[0.4em]
While each mosaic $\mathcal{M} = \{B_1, \dots, B_s\}$ is dihedrally fixed as a collective set partition, its non-focus constituent blocks $B_j$ are generally translated by a spatial root offset and may be retrograded (reflected) relative to their isolated canonical presentation in \emph{The Polygonal Burnsidean}.

Every block embedding is encoded by a compact manifest token:
\begin{equation}
\label{eqn:block-manifest-syntax}
\text{[Class]}@\text{[Root]}[\text{Chirality}]:\text{[Realised Strides]}
\end{equation}
where:
\begin{itemize}
    \item \textbf{Class}: The canonical polygon class identifier $k\text{-}i$ defined in \texttt{polyatlas} (e.g.\ $3\text{-}1$ for the isosceles triangle on $\mathbb{Z}_4$, or $3\text{-}5$ for the chiral scalene triangle on $\mathbb{Z}_8$).
    \item \textbf{Root}: The starting vertex $v_0 = \min(B_j) \in \mathbb{Z}_n$ around the circular clock face.
    \item \textbf{Chirality}: Orientation sign relative to the canonical representative:
    \begin{itemize}
        \item \textbf{$+$ (Direct)}: A pure rotational translation of the canonical polygon.
        \item \textbf{$-$ (Retrograde)}: An enantiomorphic (reflected / retrograded) copy of the canonical polygon.
    \end{itemize}
    \item \textbf{Realised Strides}: The forward cyclic strides $(s_1, \dots, s_k)$ along the perimeter of block $B_j$ starting from $v_0$.
\end{itemize}

Concatenating these block tokens across the partition yields the self-contained mosaic manifest, as illustrated in Table~\ref{tab:sample-manifests}.

\begin{table}[htbp]
\centering
\footnotesize
\setlength{\tabcolsep}{2.8pt}
\caption{\label{tab:sample-manifests}Exemplar Mosaic Orbits across $\mathbb{Z}_1..\mathbb{Z}_8$ and their Relational Manifests}
\begin{tabular}{c r c l c l l}
\toprule
Ring & ID & Shape $\lambda$ & Canonical Mosaic & $|\mathcal{O}|$ & Polygon Profile & Relational Embedding Manifest \\
\midrule
$\mathbb{Z}_1$ & 1 & $(1)$ & \texttt{/0/} & 1 & $1\text{-}1$ & \texttt{1-1@0+:1} \\
$\mathbb{Z}_2$ & 1 & $(2)$ & \texttt{/01/} & 1 & $2\text{-}1$ & \texttt{2-1@0+:11} \\
$\mathbb{Z}_2$ & 2 & $(1,1)$ & \texttt{/0/1/} & 1 & $1\text{-}1 + 1\text{-}1$ & \texttt{1-1@0+:2/1-1@1+:2} \\
$\mathbb{Z}_3$ & 1 & $(3)$ & \texttt{/012/} & 1 & $3\text{-}1$ & \texttt{3-1@0+:111} \\
$\mathbb{Z}_3$ & 2 & $(2,1)$ & \texttt{/0/12/} & 3 & $1\text{-}1 + 2\text{-}1$ & \texttt{1-1@0+:3/2-1@1+:12} \\
$\mathbb{Z}_4$ & 1 & $(4)$ & \texttt{/0123/} & 1 & $4\text{-}1$ & \texttt{4-1@0+:1111} \\
$\mathbb{Z}_4$ & 3 & $(2,2)$ & \texttt{/02/13/} & 1 & $2\text{-}2 + 2\text{-}2$ & \texttt{2-2@0+:22/2-2@1+:22} \\
$\mathbb{Z}_4$ & 4 & $(2,2)$ & \texttt{/01/23/} & 2 & $2\text{-}1 + 2\text{-}1$ & \texttt{2-1@0+:13/2-1@2+:13} \\
$\mathbb{Z}_6$ & 1 & $(6)$ & \texttt{/012345/} & 1 & $6\text{-}1$ & \texttt{6-1@0+:111111} \\
$\mathbb{Z}_6$ & 6 & $(3,3)$ & \texttt{/024/135/} & 1 & $3\text{-}2 + 3\text{-}2$ & \texttt{3-2@0+:222/3-2@1+:222} \\
$\mathbb{Z}_6$ & 8 & $(3,3)$ & \texttt{/013/245/} & 6 & $3\text{-}3 + 3\text{-}3$ & \texttt{3-3@0+:123/3-3@2-:213} \\
$\mathbb{Z}_8$ & 1 & $(8)$ & \texttt{/01234567/} & 1 & $8\text{-}1$ & \texttt{8-1@0+:11111111} \\
$\mathbb{Z}_8$ & 13 & $(4,4)$ & \texttt{/0246/1357/} & 1 & $4\text{-}6 + 4\text{-}6$ & \texttt{4-6@0+:2222/4-6@1+:2222} \\
$\mathbb{Z}_8$ & 227 & $(2^4)$ & \texttt{/01/25/37/46/} & 16 & $2\text{-}1{+}2\text{-}2{+}2\text{-}3{+}2\text{-}4$ & \texttt{2-1@0+:17/2-3@2+:35/2-4@3+:44/2-2@4+:26} \\
\bottomrule
\end{tabular}
\end{table}

In Mosaic 8 on $\mathbb{Z}_6$, for example, the two chiral triangles partition the hexagon as an enantiomorphic direct-and-reflected pair (\texttt{3-3@0+:123} and \texttt{3-3@2-:213}). In Mosaic 227 on $\mathbb{Z}_8$, the four digon chords span four distinct circular distance classes ($1, 4, 3, 2$), realising the cyclic Skolem sequence through four distinct polyatlas classes ($2\text{-}1, 2\text{-}2, 2\text{-}3, 2\text{-}4$).

The complete visual plates displaying all $509$ mosaics follow in Chapter~\ref{chap:visual-atlas}.
